\chapter{Forward Invariant Components}
\label{chp:forward_invariant_components}

In Chapter~\ref{chp:the_no_wandering_domains_theorem} we shall show that each
component of the Fatou set $F$ of $R$ has a forward image that is mapped onto
itself by some iterate $R^m$. Anticipating this, we devote this chapter to the
classification of the forward invariant components of the Fatou set of a
rational map. Although this classification logically follows
Chapter~\ref{chp:the_no_wandering_domains_theorem}, by giving it first we
emphasize that it does not depend on the more advanced material in
Chapter~\ref{chp:the_no_wandering_domains_theorem}.

% --- SECTION ------------------------------------------------------------------
% --- THE FIVE POSSIBILITIES ---------------------------------------------------
\section{The Five Possibilities}
\label{sec:the_five_possibilities}

Our objective is to give a complete analysis of the possible forward invariant
components of the Fatou set of a rational map. We shall show that such a
component can arise in exactly one of five different ways, and we begin with the
terminology used to describe these possibilities: as usual, $R$ is a rational
map of degree at least two with Fatou set $F$.

\begin{definition}[Types of forward invariant components of $F$]
    \label{def:types_of_forward_invariant_components_of_F}
    A forward invariant component $F_0$ of $F$ is:
    \begin{enumerate}[label=(\alph*)]
        \item \label{itm:attracting_component} an \emph{attracting component} if
        it contains an attracting fixed point $\zeta$ of $R$;
        \item \label{itm:superattracting_component} a \emph{super-attracting
        component} if it contains a super-attracting fixed point $\zeta$ of $R$;
        \item \label{itm:parabolic_component} a \emph{parabolic component} (or a
        \emph{Leau domain}) if there is a rationally indifferent fixed point
        $\zeta$ of $R$ on the boundary of $F_0$, and if $R^n \to \zeta$ on
        $F_0$;
        \item \label{itm:siegel_disc} a \emph{Siegel disc} if $R: F_0 \to F_0$
        is analytically conjugate to a Euclidean rotation of the unit disc onto
        itself;
        \item \label{itm:herman_ring} a \emph{Herman ring} if $R: F_0 \to F_0$
        is analytically conjugate to a Euclidean rotation of some annulus onto
        itself.
    \end{enumerate}
\end{definition}

Of course, in \ref{itm:siegel_disc} and \ref{itm:herman_ring}, the rotations
necessarily have infinite order for otherwise, some $R^n$ would be the identity
and then $\deg(R) = 1$. We know that attracting, super-attracting and parabolic
components exist, and the existence of Siegel discs was established in
Chapter~\ref{chp:periodic_points}. The proof that Herman rings exist is more
difficult, and we comment on this in
section~\ref{sec:siegel_discs_and_herman_rings}.

The central, and most important, result in this chapter is that
\ref{itm:attracting_component}-\ref{itm:herman_ring} are the only possibilities
for a forward invariant component of $F$, and this is recorded in

\begin{theorem}[There are only five possibilities for a forward invariant component of $F$]
    \label{thm:there_are_only_five_possibilities_for_a_forward_invariant_component_of_F}
    A forward invariant component of the Fatou set $F(R)$ is one of the types
    \ref{itm:attracting_component}-\ref{itm:siegel_disc} in
    Definition~\ref{def:types_of_forward_invariant_components_of_F}.
\end{theorem}
\begin{proof}
    Theorem~\ref{thm:there_are_only_five_possibilities_for_a_forward_invariant_component_of_F}
    is simply an amalgamation of several other theorems that occur later in this
    chapter, but we shall explain the general strategy of the proof now. The
    five possibilities in
    Definition~\ref{def:types_of_forward_invariant_components_of_F} can almost
    be distinguished from each other by considering those functions which can be
    expressed as the limit of some subsequence of $R^n$ in $F_0$. In
    \ref{itm:attracting_component} and \ref{itm:superattracting_component}, the
    only possible limit is the fixed point $\zeta$ and this is in $F_0$; in
    \ref{itm:parabolic_component}, the only limit function is $\zeta$, and this
    is in $\boundary{F_0}$; and finally, in \ref{itm:siegel_disc} and
    \ref{itm:herman_ring}, non-constant limit functions exist. To complete the
    description note that \ref{itm:attracting_component} and
    \ref{itm:superattracting_component} are distinguished according to whether
    $\zeta$ is, or is not, a critical point, and \ref{itm:siegel_disc} and
    \ref{itm:herman_ring} are distinguished by the connectivity of $F_0$, or by
    the existence of a fixed point in $F_0$.
\end{proof}

% --- SECTION ------------------------------------------------------------------
% --- LIMIT FUNCTIONS ----------------------------------------------------------
\section{Limit Functions}
\label{sec:limit_functions}

We now consider the class of functions which arise as locally uniform limits of
subsequences of $(R^n)_{n \in \N}$ in $F_0$. First (and regardless of whether
$F_0$ is forward invariant or not) we have

\begin{definition}[Limit function]
    A function $\varphi$ is a \emph{limit function} on a component $F_0$ of
    $F(R)$ if there is some subsequence of $(R^n)_{n \in \N}$ which converges
    locally uniformly to $\varphi$ on $F_0$. The class of limit functions on
    $F_0$ is denoted by $\mathbb{F}(F_0)$.
\end{definition}

\begin{remark}[$\mathbb{F}(F_0)$ is always non-empty and limit functions are analytic]
    As $\{R^n \mid n \geq 1\}$ is normal in $F_0$, $\mathbb{F}(F_0)$ is always
    non-empty, and each map in $\mathbb{F}(F_0)$ is analytic in $F_0$.
\end{remark}

\begin{remark}[$F_0$ forward invariant $\implies \varphi(F_0) \subset \closure{F_0}$]
    If $F_0$ is forward invariant, and if $\varphi$ is a limit function in
    $F_0$, then $\varphi(F_0)$ lies in the closure of $F_0$: in particular, if
    $\varphi$ is a constant, with value $\zeta$ say, then $\zeta \in F_0 \cup
    \boundary{F_0}$.
\end{remark}

\begin{lemma}[Constant limit functions are fixed points]
    If $F_0$ is a forward invariant component of $F$, and if there exists a
    constant limit function on $F_0$ with value $\zeta$, then $\zeta$ is a fixed
    point of $R$.
\end{lemma}
\begin{proof}
    If $z$ is in $F_0$, then so is $R(z)$, and so on some sequence of integers
    $n$ tending to $\infty$,
    \begin{equation*}
        R(\zeta) = R\left(\lim_{n \to \infty} R^n(z) \right) = \lim_{n \to \infty} R^n(R(z)) = \varphi(R(z)) = \zeta.
    \end{equation*}
\end{proof}

Of course, if the component $F_0$ is (super)attracting with fixed point $\zeta$,
say, then $R^n \to \zeta$ locally uniformly on $F_0$ and so $\mathbb{F}(F_0)$
contains a single function, namely the constant function with value $\zeta$. The
same is true (by definition) if $F_0$ is a parabolic component as in
Definition~\ref{def:types_of_forward_invariant_components_of_F}\ref{itm:parabolic_component}:
note that it is necessary to assume that $R^n \to \zeta$ in this case for a
rationally indifferent fixed point can lie on the boundary of a super-attracting
component $F$ (for example, when $R$ is a polynomial).

The next two results enable us to separate \ref{itm:attracting_component},
\ref{itm:superattracting_component} and \ref{itm:parabolic_component} from
\ref{itm:siegel_disc} and \ref{itm:herman_ring} but before giving these, we
recall 

\begin{definition}[Automorphism]
    An \emph{automorphism} of a domain $D$ is an analytic bijection of $D$ onto
    itself. The group of automorphisms of $D$ is denoted by $\automorphism{D}$.
\end{definition}

\begin{theorem}[$F_0$ forward invariant and $\mathbb{F}(F_0)$ has only constant functions $\implies \mathbb{F}(F_0) = \{\zeta\}$ fixed point and $R^n \to \zeta$ locally uniformly on $F_0$]
    \label{thm:if_all_limit_functions_are_constant_on_F0_all_limit_functions_all_constant_equals_to_the_limit_fixed_point_of_F0}
    Suppose that $F_0$ is forward invariant and that every function in
    $\mathbb{F}(F_0)$ is constant. Then $\mathbb{F}(F_0)$ contains exactly one
    function, with value $\zeta$, say, where $R(\zeta) = \zeta$ and $R^n \to
    \zeta$ locally uniformly on $F_0$.
\end{theorem}
\begin{proof}
    See \cite[163]{Beardon1991}.
\end{proof}

\begin{theorem}[$F_0$ forward invariant and $\mathbb{F}(F_0)$ has a non-constant function $\implies R \in \automorphism{F_0}$ and $\Id \in \mathbb{F}(F_0)$]
    \label{thm:if_there_is_a_nonconstant_limit_function_on_F0_all_limit_functions_are_automorphisms_and_id_is_a_limit_function}
    Suppose that $F_0$ is forward invariant, and that $\mathbb{F}(F_0)$ contains
    some non-constant function. Then $R$ is in $\automorphism{F_0}$, and the
    identity map $\Id$ is in $\mathbb{F}(F_0)$.
\end{theorem}
\begin{proof}
    See \cite[163--164]{Beardon1991}.
\end{proof}

\begin{remark}[$\varphi$ is a non-constant limit function, then $\varphi \in \automorphism{F_0}$]
    The same ideas enable one to show (in the proof of
    Theorem~\ref{thm:if_there_is_a_nonconstant_limit_function_on_F0_all_limit_functions_are_automorphisms_and_id_is_a_limit_function})
    that if $\varphi$ is a non-constant limit function, then $\varphi \in
    \automorphism{F_0}$. 
\end{remark}
\begin{proof}
    Indeed, $R \in \automorphism{F_0}$ and so has an inverse, say $S: F_0 \to
    F_0$. Now $\{S^n\}$ is a normal family in $F_0$ (because the values taken by
    $S^n$ do not lie in $J$), thus given any sequence $n_j$ with $R^{n_j} \to
    \varphi$ locally uniformly in $F_0$, we can take a subsequence $m_j$ of
    $n_j$ such that $S^{m_j}$ converges locally uniformly to some function
    $\psi$ on $F_0$. Now as $\varphi$ maps $F_0$ into itself, it follows that
    \begin{equation}
        \label{eq:psi_varphi_z_equals_z}
        \psi\varphi(z) = \lim S^{m_j}R^{m_j}(z) = z,
    \end{equation}
    and hence that $\varphi$ is injective on $F_0$. Further,
    \eqref{eq:psi_varphi_z_equals_z} shows that $\psi$ is not constant, and
    exactly as in the proof of
    Theorem~\ref{thm:if_there_is_a_nonconstant_limit_function_on_F0_all_limit_functions_are_automorphisms_and_id_is_a_limit_function},
    we see that $\psi$ maps $F_0$ into itself. With this, we can reverse the
    roles of $\varphi$ and $\psi$ in \eqref{eq:psi_varphi_z_equals_z} and deduce
    that $\varphi\psi(z) = z$ on $F_0$. It follows that $\varphi$ is a bijection
    of $F_0$ onto itself and so is in $\automorphism{F_0}$.
\end{proof}

% --- SECTION ------------------------------------------------------------------
% --- PARABOLIC DOMAINS --------------------------------------------------------
\section{Parabolic Domains}
\label{sec:parabolic_domains}

Suppose now that $F_0$ is a forward invariant component of $F(R)$ and for the
moment, regard $\boundary{F_0}$ as the ideal boundary point of the one-point
compactification of $F_0$: explicitly,
\begin{equation}
    \label{eq:ideal_boundary_point}
    R^n \to \boundary{F_0}
\end{equation}
means that for each compact subset $K$ of $F_0$, $R^n(K)$ is disjoint from $K$
for all but a finite set of $n$. If this is so, then any limit function
$\varphi$ on $F_0$ satisfies $\varphi(F_0) \subset \boundary{F_0}$, and as
$\boundary{F_0}$ has an empty interior (simply because $F_0$ is a domain), we
deduce (from the Open Mapping Theorem) that $\varphi$ is necessarily constant.
Thus, if \eqref{eq:ideal_boundary_point} holds, then by
Theorem~\ref{thm:if_all_limit_functions_are_constant_on_F0_all_limit_functions_all_constant_equals_to_the_limit_fixed_point_of_F0},
there is some fixed point $\zeta$ of $R$ in $\boundary{F_0}$ such that $R^n \to
\zeta$ locally uniformly on $F_0$. Our next result contains these conclusions
and adds one other important fact.

\begin{theorem}[$R^n \to \boundary{F_0} \implies F_0$ is a parabolic component]
    \label{thm:if_Rn_converges_on_boundary_F0_then_F0_is_a_parabolic_component}
    Suppose that $F_0$ is a forward invariant component of $F(R)$, and that $R^n
    \to \boundary{F_0}$ as $n \to \infty$. Then there is some rationally
    indifferent fixed point $\zeta$ of $R$ in $\boundary{F_0}$ such that:
    \begin{itemize}[label=(\alph*)]
        \item $R^n \to \zeta$ locally uniformly on $F_0$ as $n \to \infty$; and
        \item $R'(\zeta) = 1$.
    \end{itemize}
\end{theorem}
\begin{proof}
    See \cite[165-166]{Beardon1991}.
\end{proof}

\begin{remark}[If every limit function in $\mathbb{F}(F_0)$ is constant, $F_0$ is a (super)attracting or parabolic component]
    Suppose that $F_0$ is forward invariant and that every limit function in
    $\mathbb{F}(F_0)$ is constant. Then for some fixed point $\zeta$, $R^n \to
    \zeta$ in $F_0$ and $\zeta \in F_0 \cup \boundary{F_0}$. If $\zeta \in F_0$,
    then $F_0$ is a (super)attracting component of $F(R)$ (see the proof of
    Theorem~\ref{thm:if_all_limit_functions_are_constant_on_F0_all_limit_functions_all_constant_equals_to_the_limit_fixed_point_of_F0}),
    while if $\zeta \in \boundary{F_0}$, then by
    Theorems~\ref{thm:if_all_limit_functions_are_constant_on_F0_all_limit_functions_all_constant_equals_to_the_limit_fixed_point_of_F0}
    and
    \ref{thm:if_Rn_converges_on_boundary_F0_then_F0_is_a_parabolic_component},
    $F_0$ is a parabolic component of $F(R)$. Thus, with reference to the proof
    of
    Theorem~\ref{thm:there_are_only_five_possibilities_for_a_forward_invariant_component_of_F},
    we have shown that if every limit function in $\mathbb{F}(F_0)$ is constant,
    then $F_0$ is one of the types \ref{itm:attracting_component},
    \ref{itm:superattracting_component} and \ref{itm:parabolic_component} in
    Definition~\ref{def:types_of_forward_invariant_components_of_F}.
\end{remark}

% --- SECTION ------------------------------------------------------------------
% --- SIEGEL DISCS AND HERMAN RINGS --------------------------------------------
\section{Siegel Discs and Herman Rings}
\label{sec:siegel_discs_and_herman_rings}

In this section we prove

\begin{theorem}[$F_0$ forward invariant and $\mathbb{F}(F_0)$ has non-constant functions $\implies F_0$ is a Siegel disc or a Herman ring]
    \label{thm:if_there_are_nonconstant_limit_functions_on_F0_then_F0_is_a_Siegel_disc_or_Herman_ring}
    Suppose that $F_0$ is a forward invariant component of $F(R)$, where
    $\deg(R) \geq 2$, and that $\mathbb{F}(F_0)$ contains non-constant
    functions. Then $F_0$ is either a Siegel disc or a Herman ring.
\end{theorem}
\begin{proof}
    See \cite[167--170]{Beardon1991}.

    The idea of the proof of
    Theorem~\ref{thm:if_there_are_nonconstant_limit_functions_on_F0_then_F0_is_a_Siegel_disc_or_Herman_ring}
    is as follows. As $J$ has at least three points, the universal covering
    space of $F_0$ is the unit disc $\Delta$, and $F_0$ is conformally
    equivalent to $\Delta / \Gamma$, where $\Gamma$ is the cover group (of
    Möbius transformations) acting on $\Delta$. Now the conformal automorphisms
    $R^n$ of $F_0$ accumulate at the identity map (see
    Theorem~\ref{thm:if_there_is_a_nonconstant_limit_function_on_F0_all_limit_functions_are_automorphisms_and_id_is_a_limit_function}),
    and if we rewrite this in terms of the group $\Gamma$ and its normalizer, we
    find that $\Gamma$ is Abelian. However, the only Abelian cover groups are
    cyclic, and if we now list all possible cyclic groups and compute the
    corresponding quotient spaces, we find that $\Delta / \Gamma$, and hence
    $F_0$, is simply or doubly connected. It follows that $F_0$ is conformally
    equivalent to either a disc or an annulus and the rest of the argument is
    straightforward.
\end{proof}

\begin{remark}[Completion of the proof of Theorem~\ref{thm:there_are_only_five_possibilities_for_a_forward_invariant_component_of_F}]
    With this, we will have completed the proof of
    Theorem~\ref{thm:there_are_only_five_possibilities_for_a_forward_invariant_component_of_F}.
\end{remark}

We end this section with some brief remarks on the existence of Herman rings.
Let
\begin{equation}
    \label{eq:example_of_R_with_a_Herman_rings}
    R(z) = \lambda z^2 \left(\frac{1 + \overline{\alpha} z}{z + \alpha}\right),
\end{equation}
where $|\lambda| = 1$ and $0 < |\alpha| < 1$. If $|\alpha|$ sufficiently small,
then $R$ is a homeomorphism of the unit circle $\boundary{\Delta}$ onto itself.
Now it is known that for suitable choices of $\lambda$ and $\alpha$, $R: V \to
V$ is analytically conjugate to a rotation (necessarily of infinite order) of
$\boundary{\Delta}$ and moreover, this conjugacy extends to an $R$-invariant
neighborhood $V$ of $\boundary{\Delta}$ such that $R: V \to V$ is analytically
conjugate to a rotation of an annulus (see \cite{Arnold1965} and
\cite{Herman1979}). Clearly, $\{R^n \mid n \geq 1\}$ is normal in $V$, and if
$F_0$ is the component of $F(R)$ which contains $V$, then $\mathbb{F}(F_0)$
contains non-constant limit functions. We deduce that for these choices of
$\lambda$ and $\alpha$, $F_0$ is either a Siegel disc or a Herman ring. As both
$0$ and $\infty$ are super-attracting fixed points of $R$, $F_0$ is not simply
connected and so $F_0$ must be a Herman ring.

For an alternative construction of Herman rings, and quantitative information
about them see \cite{Shishikura1987}; this construction uses quasiconformal
maps, and we consider these a part of material in \cite{Shishikura1987} in
Chapter~8 and section~9.8 \cite{Beardon1991}.

For computer-generated illustrations of Siegel discs see
Figure~\ref{fig:figure_7_4_1}, \cite{Peitgen1986}, and for Herman rings, see
\cite{Ushiki1986}. For further information, see also references in
\cite[171]{Beardon1991}.

\begin{figure}
    \centering
    \includegraphics[width=0.8\textwidth]{beardons_book/figure_7_4_1.png}
    \caption{$z \mapsto \lambda z + z^2$, $\lambda = \exp(\pi i(\sqrt{5} -
    1))$.}
    \label{fig:figure_7_4_1}
\end{figure}

% --- SECTION ------------------------------------------------------------------
% --- CONNECTIVITY OF INVARIANT COMPONENTS -------------------------------------
\section{Connectivity of Invariant Components}
\label{sec:Connectivity_of_invariant_components}

This section contains results about the connectivity of a component $F_0$ of
$F(R)$. To prove these, we need two results about critical points, the first of
which is

\begin{theorem}[$F_0$ contains a (super)attracting fixed point $\implies F_0$ contains a critical point]
    \label{thm:F0_attracting_component_implies_F0_contains_a_critical_point}
    Let $R$ be a rational map of degree at least two. If a component $F_0$ of
    $F(R)$ contains a (super)attracting fixed point of $R$, then it also
    contains a critical point of $R$.
\end{theorem}
\begin{proof}
    See \cite[172--173]{Beardon1991}.
\end{proof}

The second result about critical points is the corresponding result for Petal
domains, namely

\begin{theorem}[$F_0$ has a rationally indifferent fixed point on the boundary $\implies F_0$ contains a critical point]
    \label{thm:F0_parabolic_component_implies_F0_contains_a_critical_point}
    Let $R$ be a rational map of degree at least two, let $\zeta$ be a
    rationally indifferent fixed point lying on the boundary of some forward
    invariant component $F_0$ of $F(R)$, and suppose that $R^n \to \zeta$
    locally uniformly in $F_0$. Then $F_0$ contains a critical point of $R$.
\end{theorem}

Although we shall not prove this result until Chapter~9 \cite{Beardon1991}
(which is devoted to the study of critical points), it is convenient (but not
necessary) to use it in the proof of the next result. For a proof, readers may
now consult Theorem~9.3.2 \cite{Beardon1991} and its proof, but for those who
would prefer not to do this, we have also included an outline of a proof of the
next result which does not use
Theorem~\ref{thm:F0_parabolic_component_implies_F0_contains_a_critical_point}.
We can now start to derive results about the connectivity of forward invariant
components of $F(R)$. First, we have

\begin{theorem}[$F_0$ forward invariant $\implies c(F_0) \in \{1, 2, \infty\}$]
    \label{thm:F0_forward_invariant_implies_cF0_is_1_2_or_infty}
    Let $R$ be a rational map of degree at least two. Then any forward invariant
    component of $F(R)$ is simply, doubly, or infinitely, connected.
\end{theorem}
\begin{proof}
    See \cite[173]{Beardon1991}.
\end{proof}

Theorem~\ref{thm:F0_forward_invariant_implies_cF0_is_1_2_or_infty} and its proof
contains the following results.

\begin{corollary}[$F$ connected $\implies c(F) \in \{1, \infty\}$]
    If $F$ is connected, then it is either simply or infinitely connected.
\end{corollary}

\begin{corollary}[$F_0$ has a fixed point $\implies c(F_0) \in \{1, \infty\}$]
    If a component of $F(R)$ contains a fixed point of $R$, then it is either
    simply or infinitely connected.
\end{corollary}

\begin{corollary}[$F_0$ doubly connected $\iff F_0$ Herman ring]
    \label{cor:F0_doubly_connected_iff_F0_Herman_ring}
    A forward invariant component of $R$ is doubly connected if and only if it
    is a Herman Ring.
\end{corollary}

\begin{remark}[$F_0$ (super)attracting or parabolic component $\implies c(F_0) \in \{1, \infty\}$]
    If a forward invariant component $F_0$ is of one of the types
    \ref{itm:attracting_component}, \ref{itm:superattracting_component} or
    \ref{itm:parabolic_component}, then according to
    Theorem~\ref{thm:F0_forward_invariant_implies_cF0_is_1_2_or_infty} and
    Corollary~\ref{cor:F0_doubly_connected_iff_F0_Herman_ring}, it must be
    simply or infinitely connected: in fact, all six possibilities can occur.
\end{remark}

We end this section with a sketch of a proof of
Theorem~\ref{thm:F0_forward_invariant_implies_cF0_is_1_2_or_infty} that does not
use
Theorem~\ref{thm:F0_parabolic_component_implies_F0_contains_a_critical_point}.
This argument is closer in spirit to the proof of
Theorem~\ref{thm:F0_attracting_component_implies_F0_contains_a_critical_point},
it requires more background knowledge, and we omit the details.

For the sketch see \cite[174--175]{Beardon1991}.